\subsection{Optimizing \aisub{s} in-sample}
\label{sec:apply_mixture}

We begin by eliciting the baseline AI's response distribution ($\hat P_0$).
We prompt \textsc{Gpt-4o} to play the basic version of the game 1,000 times without any additional instructions, setting the temperature to 1. 
Figure~\ref{fig:compare_agents_train} displays the results. 
The left panel shows the empirical PMF of the baseline AI responses (red), along with the empirical human distribution $P$ from AR's original experiment (vertical black lines). 
The baseline AI almost exclusively selects 19 shekels (\BaselineOneNineMass\%), demonstrating limited variability. 
Using the forward KL divergence as $d(\cdot,\cdot)$ with the human distribution as the reference, the distance between these distributions is $d(P,\hat P_0) =\KlHumanRawDist$.

\begin{figure}[h!]
\begin{center}
\begin{minipage}{\linewidth}
\caption{Response distributions for the basic version of the 11-20 game} 
  \vspace*{-0.15in}
\label{fig:compare_agents_train}
\centering
\includegraphics[width=.8\textwidth]{plots/basic_1120.pdf}
\end{minipage}
\end{center}
\begin{footnotesize}
\begin{singlespace}
\vspace*{-0.15in}
\emph{Notes:} This figure displays empirical PMFs for three samples playing the basic 11-20 money request game: human subjects from \citeauthor{1120Arad2012} (the vertical black lines in both panels), the off-the-shelf baseline (left panel), and responses from our selected \aisub{s} based on the weights (right panel).
The KL divergence between the human and AI distributions is reported in the upper left corner of each panel.
\end{singlespace}
\end{footnotesize}
\end{figure}

Such a poor baseline result is unsurprising.
In a related exercise, \cite{gao2024caution} also elicit responses from various LLMs playing the same 11-20 game.
Even after applying diverse prompting strategies, fine-tuning, and endowing agents with distributions of demographic traits, they similarly find that LLMs strongly index on choosing 19 shekels.
While this outcome is not inherently problematic---19 shekels lies within the support of both the Nash equilibrium and the empirical human distribution---it underscores a limitation: demographic prompts (and the other techniques they employ) alone provide little leverage in predicting strategic human reasoning.
They do not provide any reliable, flexible program for the LLM to follow, nor do they allow for heterogeneity within the simulated sample.
In contrast, the level-$k$ model implies that optimal predictions require explicitly accounting for how individuals reason about others' decisions.
And because we know from AR that there is likely a distribution of reasoning levels in their human sample, we should apply various levels of reasoning to construct a heterogeneous sample of agents.
This may, in turn, better reflect the distribution of human strategic reasoning processes.

To operationalize the level-$k$ reasoning explicitly, we construct a set of natural language \persona{s} ($\Theta_{Strategic}$) corresponding to varying levels of strategic reasoning.
These candidate \persona{s}, listed in Table~\ref{tab:persona}, specify how far ahead each \aisub{} reasons about the opponent's decisions, effectively encoding beliefs about others' strategies.

We elicit response distributions $\hat P_{\theta}$ for each candidate \persona{} $\theta \in \Theta_{Strategic}$ by prompting \textsc{Gpt-4o} 100 times per \persona{}.\footnote{We then elicit each \aisub{'s} responses using Chain-of-Thought prompting, which encourages step-by-step reasoning before producing a final answer \citep{weiCOT2024}. We implement this through two sequential prompts. \textbf{Prompt 1}: \emph{\{\texttt{11-20 game instructions}\}. Reason out a few settings according to your personality and how others might respond.} \textbf{Prompt 2}: \emph{\texttt{\{11-20 game instructions}\}. You previously had the following thoughts: \{\texttt{Response to prompt 1}\}. What amount of money would you request?}. This procedure creates the agents (and their response distributions) over which we then optimize.}
We then employ the selection method to identify the optimal mixture of \persona{s} that minimizes the absolute difference between the CDFs implied by the empirical distribution of the human responses ($P$).
This distance can be minimized using simple nonlinear programming techniques.
The weights $\mathbf{w}^*$ corresponding to the optimal mixture appear in the second column of Table~\ref{tab:persona}.\footnote{See \url{https://www.expectedparrot.com/content/6f58d11f-98cc-4de5-bb89-edcf78042d79} for the agents.}

Most mass concentrates on two \persona{s}: one that reasons between levels 1 and 3 (47\%), and another varying more broadly from levels 0 to 5 (34\%). 
The remaining weight falls on more extreme behaviors (random choices or the safest guaranteed option). 
This aligns closely with AR's findings, whose human subjects predominantly exhibited level-0 through level-3 reasoning or made random choices.

\begin{table}[h!]
\begin{center}
\caption{Proposed \aisub{} \persona{s} and resulting mixture weights from the selection method} 
\label{tab:persona}
\footnotesize
\setlength{\tabcolsep}{7pt}
\input{tables/personas.tex}
\end{center}
\begin{footnotesize}
\begin{singlespace}
\vspace*{-0.05in}
\emph{Notes:}
This table shows the set of \persona{s} $\Theta_{Strategic}$ used as input to the selection method.
The right column shows the optimized mixture weights $\mathbf{w}^*$ that minimize the absolute difference between the CDFs of the human distribution $P_s$ and the distribution of responses from the \aisub{s}.
Prepended to all the \persona{s} is: \emph{You are a human being with all the cognitive biases and heuristics that come with it}. 
We also include an explanation in the prompts for $k$-level reasoning for all \persona{s} besides the random one: \emph{A k-level thinker thinks k steps ahead. A 0-level thinker thinks 0 steps and would, therefore, just select the maximum amount that guarantees money.}
\end{singlespace}
\end{footnotesize}
\end{table}
  
Using these weights, we generate the sample $\bm{\theta}^*$ of 1,000 \aisub{s} by assigning each agent to one of the 10 \persona{s} with probability equal to its corresponding weight in Table~\ref{tab:persona}. 
The resulting empirical distribution of responses $\hat P_{\bm\theta^*}$ produced by this sample of 1,000 agents $\bm{\theta}^*$ appears in the right panel of Figure~\ref{fig:compare_agents_train} (blue).
The improvement over the baseline AI is substantial: $d(P, \hat P_{\bm\theta^*}) = \KlHumanOptDist$ is \KlPercDiffBasic\% smaller than $d(P, \hat P_0) = \KlHumanRawDist$, demonstrating a strong in-sample fit.
  